\section{Quản lý đồ án và phụ lục}
\label{sec:contribution}

\subsection{Kế hoạch triển khai}
Để bảo đảm tiến độ đồng đều và tránh tình trạng chờ đợi dây chuyền giữa bốn thành
viên trong nhóm, dự án được tổ chức theo phương pháp thiết kế đặc tả giao tiếp
(\textit{API contract}) trước khi cài đặt chi tiết, chia làm ba giai đoạn nối tiếp
nhau rõ ràng. Trong giai đoạn khởi đầu, toàn nhóm cùng phân tích đặc tả bài toán,
thống nhất mô hình dữ liệu cốt lõi, quy ước định dạng mệnh đề số nguyên và ranh
giới truyền nhận dữ liệu bất biến giữa bộ điều hành trò chơi với tác tử suy diễn.
Nhờ bộ khung giao tiếp chuẩn hóa này, sang giai đoạn phát triển trọng tâm, cả bốn
phân hệ lớn --- bao gồm môi trường quản lý trạng thái ẩn kèm bộ mã hóa CNF, bộ
giải thỏa mãn mệnh đề DPLL tự cài đặt, tác tử suy diễn phản chứng kèm bộ nhớ đệm
chứng nhân, và tầng dịch vụ API kết nối trực tiếp kèm bộ công cụ đo kiểm hiệu năng
(\textit{benchmark}) trên tiến trình độc lập --- đều được lập trình và kiểm thử đơn
vị song song trên các nhánh làm việc riêng biệt. Ở giai đoạn tích hợp và hoàn
thiện, các nhánh thành phần được rà soát chéo và hợp nhất thông qua hệ thống kiểm
thử hồi quy tự động gồm $290$ ca kiểm thử, trước khi nâng cấp không gian giao diện
ba chiều, tiến hành $42$ lượt đo đạc thực nghiệm chính thức và biên soạn báo cáo
tổng kết.

\subsection{Phân công nhiệm vụ}
Bảng \ref{tab:team_contributions} trình bày chi tiết khối lượng công việc thực tế
do từng thành viên đảm nhận (đối chiếu trung thực với lịch sử đóng góp trên kho
lưu trữ chung của nhóm) cùng mức độ hoàn thành nhiệm vụ được giao. Theo quy chế
đánh giá của học phần, tỷ lệ hoàn thành $100\%$ ở cả bốn thành viên xác nhận mỗi
cá nhân đều đã thực hiện đầy đủ, đúng hạn và đạt yêu cầu chất lượng phần việc
chuyên môn được phân công.

\begin{table}[htbp]
\centering
\caption{%
    Bảng phân công nhiệm vụ chuyên môn và tỷ lệ hoàn thành của các thành viên
    trong nhóm}
\label{tab:team_contributions}
\setstretch{1.08}
\scriptsize
\renewcommand{\arraystretch}{1.2}
\begin{tabularx}{\textwidth}{p{3.4cm} l >{\raggedright\arraybackslash}X c}
\toprule
\textbf{Họ và Tên} & \textbf{MSSV} & \textbf{Nội dung công việc đảm nhiệm thực tế} &
\textbf{Hoàn thành} \\
\midrule
\textbf{Nguyễn Phúc Khang}\newline \textit{(Trưởng nhóm)} & 24120068 &
Định hình kiến trúc phân tầng tổng thể, thiết kế mô hình dữ liệu cốt lõi và điều
phối rà soát hợp nhất các nhánh phát triển của nhóm. Xây dựng bộ quản lý trạng thái
ẩn, quy trình kiểm định tính trung thực lời khai, kho cơ sở dữ liệu $14$ màn chơi
chuẩn và bộ đánh giá ngữ nghĩa đối chứng độc lập. Nghiên cứu và hiện thực bộ mã
hóa tự động sang dạng chuẩn hội cho sáu họ manh mối cốt lõi cùng hai mở rộng giao
vùng hình học và mạch chẵn lẻ biến phụ Tseitin tuyến tính. Thiết kế không gian bàn
điều tra ba chiều \textit{Midnight Archive} với hiệu ứng ánh sáng vật lý, hiệu ứng
shader mưa đêm, con dấu đóng trực tiếp lên thẻ và tổng hợp biên tập báo cáo tổng
kết đồ án. & \textbf{100\%} \\
\midrule
\textbf{Phạm Minh Trọng} & 24120151 &
Thiết kế và hiện thực tầng dịch vụ API quản lý định danh phiên làm việc, đồng bộ số
hiệu phiên bản tương tranh lạc quan và kết nối luồng kiểm chứng phán quyết, gợi ý
cùng giải tự động từng bước. Xây dựng bộ công cụ đo kiểm hiệu năng chạy trên tiến
trình hệ điều hành độc lập với cơ chế áp đặt giới hạn thời gian cứng, thu thập chỉ
số đo đạc trên $14$ màn chơi và tự động tổng hợp dữ liệu thực nghiệm. Phát triển
hệ thống kiểm thử tích hợp đầu--cuối và rà soát đối chiếu số liệu thống kê trong
báo cáo. & \textbf{100\%} \\
\midrule
\textbf{Trần Nguyễn Anh Khoa} & 24120075 &
Nghiên cứu lý thuyết và hiện thực từ đầu bộ giải thỏa mãn mệnh đề DPLL tất định
thuần túy không sử dụng thư viện giải ngoài. Xây dựng cơ chế tiếp nhận giả định
tạm thời không làm biến đổi cơ sở tri thức, thủ tục lan truyền mệnh đề đơn vị lặp
đến điểm bất động, phát hiện xung đột mệnh đề rỗng sớm, chiến lược chọn biến rẽ
nhánh tất định kèm sao chép cô lập trạng thái quay lui và bộ đếm thống kê thực thi
chi tiết. Thiết kế bộ kiểm thử chuyên sâu cho bộ giải trên các công thức biên và
công thức bất khả thỏa lớn. & \textbf{100\%} \\
\midrule
\textbf{Bùi Nhật Bảo} & 24120167 &
Nghiên cứu và hiện thực tác tử suy diễn logic dựa trên nguyên lý chứng minh phản
chứng bốn trạng thái tuân thủ chặt chẽ quy tắc không suy đoán thiếu căn cứ. Đề xuất
ý tưởng, chứng minh tính đúng đắn và cài đặt cơ chế bộ nhớ đệm mô hình chứng nhân
theo từng bước suy diễn nhằm cắt giảm các lượt giải dư thừa. Xây dựng chu trình
giải tự động theo thứ tự ưu tiên hàng có ghi vết suy diễn đối soát từng bước, thuật
toán sinh gợi ý an toàn từ dữ kiện công khai và quy trình kiểm định nghiệm duy nhất
sử dụng mệnh đề chặn. & \textbf{100\%} \\
\bottomrule
\end{tabularx}
\end{table}

\subsection{Tài nguyên dự án}
Toàn bộ mã nguồn của hệ thống, tập dữ liệu $14$ màn chơi chuẩn, bộ $290$ kịch bản
kiểm thử tự động, công cụ tái lập thực nghiệm cùng video minh họa toàn bộ luồng
hoạt động trên không gian bàn điều tra ba chiều (quay ở độ phân giải cao
$1600 \times 900$ tốc độ $30$ khung hình/giây, trình diễn đầy đủ thao tác quan sát
cảnh ba chiều, lọc vùng manh mối, phản hồi đóng dấu từ chối trực tiếp lên mặt thẻ
khi kết luận sớm, sử dụng gợi ý và giải tự động tới khi khép lại vụ án) được lưu
trữ công khai tại kho mã nguồn trực tuyến của nhóm:
\url{https://github.com/khang1108/griductive}.

\subsection{Tuyên bố sử dụng trí tuệ nhân tạo}
\label{subsec:ai_appendix}
Tuân thủ nguyên tắc liêm chính học thuật và yêu cầu minh bạch tại Mục 4.6 của đề
cương đồ án, nhóm xin khai báo cụ thể phạm vi ứng dụng các mô hình ngôn ngữ lớn
trong quá trình thực hiện dự án. Trong suốt chu kỳ phát triển, công cụ trợ lý AI
chỉ được sử dụng như một phương tiện hỗ trợ kỹ thuật ở các khâu phụ trợ gồm: tham
khảo cách tổ chức ánh sáng vật lý, vật liệu bề mặt gỗ và kỹ thuật viết chương trình
shader cửa sổ mưa đêm trong môi trường đồ họa ba chiều trên trình duyệt; gợi ý mở
rộng thêm các kịch bản kiểm thử biên tự động để nâng cao độ bao phủ kiểm thử phần
mềm; đồng thời hỗ trợ rà soát lỗi chính tả, căn chỉnh bố cục bảng biểu và sơ đồ
khối trong quá trình chế bản tài liệu \LaTeX.

Để làm rõ cách thức tương tác có kiểm soát của nhóm với công cụ hỗ trợ, dưới đây
là ba yêu cầu truy vấn (\textit{prompt}) tiêu biểu đại diện cho các tình huống tra
cứu kỹ thuật thực tế trong dự án:
\begin{softbox}
\textbf{Prompt mẫu 1 (Kiểm chứng chéo tính tương đương ngữ nghĩa của bộ mã hóa):}
\textit{``Hãy đề xuất kịch bản kiểm thử vét cạn toàn bộ phép gán chân trị trên lưới
$3 \times 3$ và $4 \times 4$ để đối chiếu tập cấu hình thỏa mãn theo định nghĩa ngữ
nghĩa gốc với kết quả khả thỏa của công thức chuẩn hội sau khi mã hóa, đặc biệt chú
ý các trường hợp biên của ràng buộc chẵn lẻ dùng biến phụ Tseitin và vùng giao hình
học.''}\\[6pt]
\textbf{Prompt mẫu 2 (Thiết kế hiệu ứng phản hồi trực quan ba chiều tối giản):}
\textit{``Làm thế nào để thiết kế hoạt ảnh một con dấu cao su nghiêng dập trực
tiếp lên bề mặt thẻ nghi phạm trong không gian ba chiều rồi từ từ mờ dần khi người
chơi đưa ra phán quyết chưa đủ cơ sở chứng minh, nhằm thay thế các hộp thoại thông
báo gây rối mắt trên màn hình?''}\\[6pt]
\textbf{Prompt mẫu 3 (Chế bản học thuật trong \LaTeX):} \textit{``Hướng dẫn cách
căn chỉnh sơ đồ kiến trúc phân tầng bằng thư viện vẽ vector trong \LaTeX{} để thể
hiện rõ luồng dữ liệu một chiều chỉ trích xuất trạng thái công khai từ bộ điều hành
trò chơi sang tầng tác tử suy diễn.''}
\end{softbox}

Nhóm xin cam kết và chịu trách nhiệm rằng toàn bộ phần mô hình hóa toán học biến
mệnh đề, các phép biến đổi đại số tương đương sang dạng chuẩn hội, thuật toán giải
thỏa mãn mệnh đề DPLL tự cài đặt, định lý phản chứng cùng cơ chế bộ nhớ đệm mô
hình chứng nhân của tác tử suy diễn, cũng như toàn bộ số liệu thực nghiệm thu thập
từ $42$ lượt đo đạc độc lập và nội dung phân tích học thuật trong báo cáo này đều
do các thành viên trong nhóm trực tiếp nghiên cứu, thiết kế, kiểm chứng chéo và
biên soạn.
